\documentclass[12pt]{article}
\usepackage[answers]{signals-style}   % shared style: styles/signals-style.sty

\begin{document}

\tutorialpage{01}{September 2026}

\begin{enumerate}
\item \mbox{}

% minipage + \captionsetup{type=figure} keeps the figure inside this
% item instead of letting it float away
\par\noindent
\begin{minipage}{\linewidth}
\centering
\captionsetup{type=figure}
    % Common style for all six plots (local to this figure)
    \pgfplotsset{
        sigplot/.style={
            scale only axis,
            width = 0.8\linewidth,
            height = 2.6cm,
            axis lines = middle,
            samples = 200,
            grid = both,
            tick label style = {font=\scriptsize, fill=white, inner sep=1pt},
            xlabel style = {at={(ticklabel* cs:1)}, anchor=west},
            ylabel style = {at={(ticklabel* cs:1)}, anchor=south},
            enlarge x limits = {abs=0.2},
            enlarge y limits = {abs=0.2},
        }
    }
    \begin{subfigure}{0.32\textwidth}
    \centering
    \begin{tikzpicture}
        \begin{axis}[
        sigplot,
        xlabel = $t$,
        ylabel = $x(t)$,
        xmin = -1.5, xmax = 1.5,
        ymin = -1, ymax = 1,
        xtick distance = 1,
        minor x tick num = 3,
        ytick = {-1,0,1},
        ]
        \addplot[domain=-1.5:1.5, thick]{sin(deg(2*pi*x))};
        \end{axis}
    \end{tikzpicture}
    \caption{$x(t) = \sin(2\pi t)$}
    \label{fig:q1a}
    \end{subfigure}
\hfill
    \begin{subfigure}{0.32\textwidth}
    \centering
    \begin{tikzpicture}
        \begin{axis}[
        sigplot,
        xlabel = $t$,
        ylabel = $x(t)$,
        xmin = -12, xmax = 12,
        ymin = -1, ymax = 1,
        xtick distance = 4,
        minor x tick num = 1,
        ytick = {-1,0,1},
        enlarge x limits = {abs=1.5},
        ]
        \addplot[domain=-12:12, thick]{cos(deg(pi*x/4))};
        \end{axis}
    \end{tikzpicture}
    \caption{$x(t) = \cos(\pi t/4)$}
    \label{fig:q1b}
    \end{subfigure}
\hfill
    \begin{subfigure}{0.32\textwidth}
    \centering
    \begin{tikzpicture}
        \begin{axis}[
        sigplot,
        xlabel = $t$,
        ylabel = $x(t)$,
        xmin = -4, xmax = 6,
        ymin = -1, ymax = 1,
        xtick distance = 2,
        minor x tick num = 1,
        ytick = {-1,0,1},
        enlarge x limits = {abs=0.6},
        ]
        \addplot[domain=-4:6, thick]{sin(deg(pi*x/2 - pi/6))};
        % zero crossing shifted to t = 1/3
        \addplot[dashed, gray] coordinates{(1/3,-1) (1/3,1)};
        % y-axis intercept: x(0) = sin(-pi/6) = -1/2
        \addplot[only marks, mark=*, mark size=1.5pt] coordinates{(0,-0.5)};
        \node[anchor=south east, font=\scriptsize, inner sep=1pt] at (axis cs:0,-0.5) {$-\tfrac{1}{2}$};
        \end{axis}
    \end{tikzpicture}
    \caption{$x(t) = \sin(\pi t/2 - \pi/6)$}
    \label{fig:q1c}
    \end{subfigure}

    \vspace{1em}

    \begin{subfigure}{0.32\textwidth}
    \centering
    \begin{tikzpicture}
        \begin{axis}[
        sigplot,
        xlabel = $t$,
        ylabel = $x(t)$,
        xmin = -1, xmax = 8,
        ymin = 0, ymax = 1.2,
        xtick = {5},
        minor xtick = {-1,...,8},
        ytick = {0,1},
        enlarge x limits = {abs=0.5},
        ]
        \addplot[domain=-1:5, thick]{0};
        \addplot[thick] coordinates{(5,0) (5,1)};
        \addplot[domain=5:8, thick]{1};
        \end{axis}
    \end{tikzpicture}
    \caption{$x(t) = u(t-5)$}
    \label{fig:q1d}
    \end{subfigure}
\hfill
    \begin{subfigure}{0.32\textwidth}
    \centering
    \begin{tikzpicture}
        \begin{axis}[
        sigplot,
        xlabel = $t$,
        ylabel = $x(t)$,
        xmin = -1, xmax = 4,
        ymin = 0, ymax = 1.2,
        xtick distance = 1,
        ytick = {0,1},
        ]
        \addplot[domain=-1:0, thick]{0};
        \addplot[thick] coordinates{(0,0) (0,1)};
        \addplot[domain=0:4, thick]{exp(-x)};
        \end{axis}
    \end{tikzpicture}
    \caption{$x(t) = e^{-t}u(t)$}
    \label{fig:q1e}
    \end{subfigure}
\hfill
    \begin{subfigure}{0.32\textwidth}
    \centering
    \begin{tikzpicture}
        \begin{axis}[
        sigplot,
        xlabel = $n$,
        ylabel = $x{[n]}$,
        xmin = -12, xmax = 12,
        ymin = -1, ymax = 1,
        xtick distance = 6,
        minor x tick num = 1,
        ytick = {-1,0,1},
        enlarge x limits = {abs=1.5},
        ]
        \addplot[ycomb, thick, mark=*, mark size=1.2pt,
                 domain=-12:12, samples=25]{sin(deg(pi*x/6))};
        \end{axis}
    \end{tikzpicture}
    \caption{$x[n] = \sin(\pi n/6)$}
    \label{fig:q1f}
    \end{subfigure}
\caption{Sketches of the signals in Question 1.}
\label{fig:q1}
\end{minipage}

% ---------------- Q2 ----------------
\item \mbox{}

\par\noindent
\begin{minipage}{\linewidth}
\centering
\captionsetup{type=figure}
    % Same common style as Question 1 (local to this figure)
    \pgfplotsset{
        sigplot/.style={
            scale only axis,
            width = 0.8\linewidth,
            height = 2.6cm,
            axis lines = middle,
            grid = both,
            tick label style = {font=\scriptsize, fill=white, inner sep=1pt},
            xlabel style = {at={(ticklabel* cs:1)}, anchor=west},
            ylabel style = {at={(ticklabel* cs:1)}, anchor=south},
            enlarge x limits = {abs=0.2},
            enlarge y limits = {abs=0.2},
        }
    }
    \begin{subfigure}{0.32\textwidth}
    \centering
    \begin{tikzpicture}
        \begin{axis}[
        sigplot,
        xlabel = $t$,
        ylabel = $x(t)+1$,
        xmin = -1, xmax = 3,
        ymin = 0, ymax = 2,
        xtick distance = 1,
        ytick = {0,1,2},
        ]
        \addplot[thick] coordinates{(-1,1) (0,1) (1,2) (2,2) (2,1) (3,1)};
        \end{axis}
    \end{tikzpicture}
    \caption{$x(t)+1$}
    \label{fig:q2a}
    \end{subfigure}
\hfill
    \begin{subfigure}{0.32\textwidth}
    \centering
    \begin{tikzpicture}
        \begin{axis}[
        sigplot,
        xlabel = $t$,
        ylabel = $2x(t)-1$,
        xmin = -1, xmax = 3,
        ymin = -1, ymax = 1,
        xtick distance = 1,
        ytick = {-1,0,1},
        ]
        \addplot[thick] coordinates{(-1,-1) (0,-1) (1,1) (2,1) (2,-1) (3,-1)};
        \end{axis}
    \end{tikzpicture}
    \caption{$2x(t)-1$}
    \label{fig:q2b}
    \end{subfigure}
\hfill
    \begin{subfigure}{0.32\textwidth}
    \centering
    \begin{tikzpicture}
        \begin{axis}[
        sigplot,
        xlabel = $t$,
        ylabel = $-x(t)$,
        xmin = -1, xmax = 3,
        ymin = -1, ymax = 0.5,
        xtick distance = 1,
        ytick = {-1,0},
        ]
        \addplot[thick] coordinates{(-1,0) (0,0) (1,-1) (2,-1) (2,0) (3,0)};
        \end{axis}
    \end{tikzpicture}
    \caption{$-x(t)$}
    \label{fig:q2c}
    \end{subfigure}

    \vspace{1em}

    \begin{subfigure}{0.32\textwidth}
    \centering
    \begin{tikzpicture}
        \begin{axis}[
        sigplot,
        xlabel = $t$,
        ylabel = $-2x(t)+1$,
        xmin = -1, xmax = 3,
        ymin = -1, ymax = 1,
        xtick distance = 1,
        ytick = {-1,0,1},
        ]
        \addplot[thick] coordinates{(-1,1) (0,1) (1,-1) (2,-1) (2,1) (3,1)};
        \end{axis}
    \end{tikzpicture}
    \caption{$-2x(t)+1$}
    \label{fig:q2d}
    \end{subfigure}
\hfill
    \begin{subfigure}{0.32\textwidth}
    \centering
    \begin{tikzpicture}
        \begin{axis}[
        sigplot,
        xlabel = $t$,
        ylabel = $x(t-1)$,
        xmin = -1, xmax = 4,
        ymin = 0, ymax = 1.2,
        xtick distance = 1,
        ytick = {0,1},
        ]
        \addplot[thick] coordinates{(-1,0) (1,0) (2,1) (3,1) (3,0) (4,0)};
        \end{axis}
    \end{tikzpicture}
    \caption{$x(t-1)$}
    \label{fig:q2e}
    \end{subfigure}
\hfill
    \begin{subfigure}{0.32\textwidth}
    \centering
    \begin{tikzpicture}
        \begin{axis}[
        sigplot,
        xlabel = $t$,
        ylabel = $x(-t)$,
        xmin = -3, xmax = 1,
        ymin = 0, ymax = 1.2,
        xtick distance = 1,
        ytick = {0,1},
        ]
        \addplot[thick] coordinates{(-3,0) (-2,0) (-2,1) (-1,1) (0,0) (1,0)};
        \end{axis}
    \end{tikzpicture}
    \caption{$x(-t)$}
    \label{fig:q2f}
    \end{subfigure}

    \vspace{1em}

    \begin{subfigure}{0.32\textwidth}
    \centering
    \begin{tikzpicture}
        \begin{axis}[
        sigplot,
        xlabel = $t$,
        ylabel = $x(-t+1)$,
        xmin = -2, xmax = 2,
        ymin = 0, ymax = 1.2,
        xtick distance = 1,
        ytick = {0,1},
        ]
        \addplot[thick] coordinates{(-2,0) (-1,0) (-1,1) (0,1) (1,0) (2,0)};
        \end{axis}
    \end{tikzpicture}
    \caption{$x(-t+1)$}
    \label{fig:q2g}
    \end{subfigure}
\hfill
    \begin{subfigure}{0.32\textwidth}
    \centering
    \begin{tikzpicture}
        \begin{axis}[
        sigplot,
        xlabel = $t$,
        ylabel = $x(2t)$,
        xmin = -1, xmax = 2,
        ymin = 0, ymax = 1.2,
        xtick distance = 0.5,
        ytick = {0,1},
        ]
        \addplot[thick] coordinates{(-1,0) (0,0) (0.5,1) (1,1) (1,0) (2,0)};
        \end{axis}
    \end{tikzpicture}
    \caption{$x(2t)$}
    \label{fig:q2h}
    \end{subfigure}
\hfill
    \begin{subfigure}{0.32\textwidth}
    \centering
    \begin{tikzpicture}
        \begin{axis}[
        sigplot,
        xlabel = $t$,
        ylabel = $x(t/2+1)$,
        xmin = -3, xmax = 3,
        ymin = 0, ymax = 1.2,
        xtick distance = 1,
        ytick = {0,1},
        ]
        \addplot[thick] coordinates{(-3,0) (-2,0) (0,1) (2,1) (2,0) (3,0)};
        \end{axis}
    \end{tikzpicture}
    \caption{$x(t/2+1)$}
    \label{fig:q2i}
    \end{subfigure}
\caption{Sketches of the signals in Question 2.}
\label{fig:q2}
\end{minipage}

% ---------------- Q3 ----------------
\item \mbox{}

\par\noindent
\begin{minipage}{\linewidth}
\centering
\captionsetup{type=figure}
    % Same common style, with stem plots for discrete-time signals
    \pgfplotsset{
        sigplot/.style={
            scale only axis,
            width = 0.8\linewidth,
            height = 2.6cm,
            axis lines = middle,
            grid = both,
            tick label style = {font=\scriptsize, fill=white, inner sep=1pt},
            xlabel style = {at={(ticklabel* cs:1)}, anchor=west},
            ylabel style = {at={(ticklabel* cs:1)}, anchor=south},
            enlarge x limits = {abs=0.6},
            enlarge y limits = {abs=0.4},
            xtick distance = 1,
        },
        stem/.style={ycomb, thick, mark=*, mark size=1.2pt},
    }
    \begin{subfigure}{0.32\textwidth}
    \centering
    \begin{tikzpicture}
        \begin{axis}[
        sigplot,
        xlabel = $n$,
        ylabel = ${2x[n]-1}$,
        xmin = -0.5, xmax = 5.5,
        xticklabel style = {fill=none},  % don't hide negative stems
        ymin = -1, ymax = 5,
        ytick = {-1,1,3,5},
        ]
        \addplot[stem] coordinates{(0,1) (1,3) (2,5) (3,3) (4,1) (5,-1)};
        \end{axis}
    \end{tikzpicture}
    \caption{$2x[n]-1$}
    \label{fig:q3a}
    \end{subfigure}
\hfill
    \begin{subfigure}{0.32\textwidth}
    \centering
    \begin{tikzpicture}
        \begin{axis}[
        sigplot,
        xlabel = $n$,
        ylabel = ${-x[n]+1}$,
        xmin = 0, xmax = 5,
        xticklabel style = {fill=none},  % don't hide negative stems
        ymin = -2, ymax = 1,
        ytick = {-2,-1,0,1},
        ]
        \addplot[stem] coordinates{(0,0) (1,-1) (2,-2) (3,-1) (4,0) (5,1)};
        \end{axis}
    \end{tikzpicture}
    \caption{$-x[n]+1$}
    \label{fig:q3b}
    \end{subfigure}
\hfill
    \begin{subfigure}{0.32\textwidth}
    \centering
    \begin{tikzpicture}
        \begin{axis}[
        sigplot,
        xlabel = $n$,
        ylabel = ${x[n+1]}$,
        xmin = -1, xmax = 4,
        ymin = 0, ymax = 3,
        ytick = {0,1,2,3},
        ]
        \addplot[stem] coordinates{(-1,1) (0,2) (1,3) (2,2) (3,1) (4,0)};
        \end{axis}
    \end{tikzpicture}
    \caption{$x[n+1]$}
    \label{fig:q3c}
    \end{subfigure}

    \vspace{1em}

    \begin{subfigure}{0.32\textwidth}
    \centering
    \begin{tikzpicture}
        \begin{axis}[
        sigplot,
        xlabel = $n$,
        ylabel = ${x[-n+2]}$,
        xmin = -3, xmax = 2,
        ymin = 0, ymax = 3,
        ytick = {0,1,2,3},
        ]
        \addplot[stem] coordinates{(-3,0) (-2,1) (-1,2) (0,3) (1,2) (2,1)};
        \end{axis}
    \end{tikzpicture}
    \caption{$x[-n+2]$}
    \label{fig:q3d}
    \end{subfigure}
\hfill
    \begin{subfigure}{0.32\textwidth}
    \centering
    \begin{tikzpicture}
        \begin{axis}[
        sigplot,
        xlabel = $n$,
        ylabel = ${x[2n+1]}$,
        xmin = 0, xmax = 2,
        ymin = 0, ymax = 3,
        ytick = {0,1,2,3},
        ]
        % only x[1], x[3], x[5] are picked out
        \addplot[stem] coordinates{(0,2) (1,2) (2,0)};
        \end{axis}
    \end{tikzpicture}
    \caption{$x[2n+1]$}
    \label{fig:q3e}
    \end{subfigure}
\hfill
    \begin{subfigure}{0.32\textwidth}
    \centering
    \begin{tikzpicture}
        \begin{axis}[
        sigplot,
        xlabel = $n$,
        ylabel = ${x[n/2]}$,
        xmin = 0, xmax = 10,
        ymin = 0, ymax = 3,
        xtick distance = 2,
        minor x tick num = 1,
        ytick = {0,1,2,3},
        ]
        % defined only for even n
        \addplot[stem] coordinates{(0,1) (2,2) (4,3) (6,2) (8,1) (10,0)};
        \end{axis}
    \end{tikzpicture}
    \caption{$x[n/2]$}
    \label{fig:q3f}
    \end{subfigure}
\caption{Sketches of the signals in Question 3.}
\label{fig:q3}
\end{minipage}

% ---------------- Q4 ----------------
\item A continuous-time signal $\sin(\omega t + \theta)$ is periodic with $T = 2\pi/\omega$.
A discrete-time signal $\sin(\Omega n)$ is periodic only if $\Omega N = 2\pi k$ for some integers $N, k$, i.e.\ only if $\Omega/2\pi$ is rational.
\begin{enumerate}
    \item $\omega = \pi$, so $T = 2\pi/\pi = 2$.
    \textbf{Periodic, $T = 2$.}

    \item The phase shift $\pi/4$ does not change the period; $\omega = \pi$, so $T = 2$.
    \textbf{Periodic, $T = 2$.}

    \item $\sin(\pi t/4)$ alone has period $8$, but the envelope $e^t$ keeps growing:
    \[x(t+8) = e^{t+8}\sin(\pi t/4) = e^{8}\,x(t) \neq x(t).\]
    No $T>0$ gives $x(t+T) = x(t)$ for all $t$. \textbf{Not periodic.}

    \item $\Omega = \pi/4$. We need $\frac{\pi}{4}N = 2\pi k \Rightarrow N = 8k$.
    The smallest positive integer is $N = 8$ (with $k=1$). \textbf{Periodic, $N = 8$.}

    \item $\Omega = 2$. We need $2N = 2\pi k \Rightarrow N = \pi k$, which is never an integer for $k \neq 0$
    ($\Omega/2\pi = 1/\pi$ is irrational). \textbf{Not periodic.}
\end{enumerate}

\end{enumerate}

\end{document}
